\begin{table}[!htbp]
\centering
\caption{Observed fund trading and return per unit of measured demand}\label{tab:demand_main}
\small
\setlength{\tabcolsep}{5pt}
\begin{tabular}{llrrrrl}
\toprule
Type & Period & Events & Return & Demand & $M$ & 95\% ratio set \\
\midrule
Addition & 2007--11 & 8 & 8.78 & 0.15 & 57.87 & [-9.55, 471.88] \\
Addition & 2012--16 & 27 & -1.94 & 0.52 & -3.76 & [-15.09, 4.24] \\
Addition & 2017--21 & 58 & 3.29 & 1.11 & 2.97 & [-0.14, 6.30] \\
Addition & 2022--26 & 78 & 2.28 & 0.97 & 2.35 & [-0.21, 4.73] \\
Addition & 2017--26 & 136 & 2.71 & 1.03 & 2.63 & [0.91, 4.37] \\
Deletion & 2007--11 & 7 & -0.19 & 0.26 & -0.72 & All real values \\
Deletion & 2012--16 & 25 & 2.07 & 0.48 & 4.36 & [-17.73, 30.17] \\
Deletion & 2017--21 & 40 & 0.20 & 0.99 & 0.20 & [-4.05, 4.89] \\
Deletion & 2022--26 & 68 & 1.66 & 1.18 & 1.41 & [-0.87, 2.83] \\
Deletion & 2017--26 & 108 & 1.12 & 1.11 & 1.01 & [-0.88, 2.61] \\
\bottomrule
\end{tabular}
\begin{minipage}{\linewidth}\footnotesize\vspace{5pt}
Return is sign-adjusted CAR from AD-1 through ED+1, in percent. Demand is the sign-adjusted broad benchmark-linked holdings change as a percentage of shares outstanding. For deletions, negative returns and sales have positive signs. Each ratio divides the two means on identical events retained after the share-adjustment screen. Joint ratio uncertainty is described in Appendix~\ref{app:inference}; it can be unbounded when the denominator is weak. These are descriptive ratios, not identified elasticities.
\end{minipage}
\end{table}
